\hypersetup{pageanchor=false}

\begin{titlepage}
  \centering
  \vspace*{0.18\textheight}
  {\Huge\bfseries Navier--Stokes 方程\par}
  \vspace{0.8em}
  {\LARGE 理论与数值分析\par}
  \vspace{3em}
  {\Large ROGER TEMAM\par}
  \vfill
  {\large AMS CHELSEA PUBLISHING\par}
  \vspace{0.5em}
  {American Mathematical Society\par}
  {Providence, Rhode Island\par}
\end{titlepage}

\clearpage
\thispagestyle{empty}
\null
\clearpage

\begin{titlepage}
  \centering
  \vspace*{0.16\textheight}
  {\Huge\bfseries Navier--Stokes 方程\par}
  \vspace{0.8em}
  {\LARGE 理论与数值分析\par}
  \vspace{2.5em}
  {\large 作者\par}
  \vspace{0.5em}
  {\Large ROGER TEMAM\par}
  \vfill
  {AMS CHELSEA PUBLISHING\par}
  {American Mathematical Society, Providence, Rhode Island\par}
\end{titlepage}

\clearpage
\hypersetup{pageanchor=true}
\pagenumbering{roman}
\thispagestyle{empty}
\section*{出版信息}

2000 Mathematics Subject Classification: Primary 76--02; Secondary 65--02, 76D05, 35Q30.

Library of Congress Cataloging-in-Publication Data: Temam, Roger. \emph{Navier--Stokes equations: theory and numerical analysis}. 本书最初于 1977 年由 North-Holland 在 Amsterdam 和 New York 出版, 含参考文献与索引. ISBN 0-8218-2737-5 (alk. paper). 分类号 QA374 .T44 2001; 532'.0527'01515353--dc21; 00-067641.

个人读者以及代表个人读者的非营利图书馆, 可以在合理使用范围内复制本书材料, 例如为教学或研究复制一章. 书评可以引用短篇幅内容, 但须按惯例注明来源. 以再出版, 系统复制或多份复制为目的使用本书任何材料, 必须取得 American Mathematical Society 的许可. 许可申请应寄送至 Assistant to the Publisher, American Mathematical Society, P. O. Box 6248, Providence, Rhode Island 02940-6248, 或发送电子邮件至 \texttt{reprint-permission@ams.org}.

Copyright \textcopyright\ 1984, currently held by the American Mathematical Society. American Mathematical Society 于 2001 年重印并订正. Printed in the United States of America. 本书所用纸张无酸, 符合关于永久保存性和耐久性的规范. AMS 主页: \url{http://www.ams.org/}.

\clearpage
\thispagestyle{empty}
\phantomsection
\label{front:dedication}
\vspace*{0.36\textheight}
\begin{center}
  谨以崇高敬意献给 Jean Leray.
\end{center}
\clearpage

\phantomsection
\label{front:contents}
\addcontentsline{toc}{section}{目录}
\tableofcontents
\clearpage

\section*{AMS Chelsea 版序言}
\phantomsection
\label{front:preface-ams-chelsea}
\addcontentsline{toc}{section}{AMS Chelsea 版序言}

本版重现了 North-Holland 于 1977 年首次出版的本书. 在形式上, 全书由 AMS 重新排版; 在内容上, 除少量编辑性修改外, 与 1984 年出版的第三次修订版相同. 如果今天重新撰写, 本书很可能会在若干方面有所不同. 另一方面, 在新版中引入修改需要大量工作, 结果未必可靠, 而且很可能带来新的错误. 因此, 我们决定按上一版的原貌重印本书.

本书新增的材料是 \MBUnresolvedReference{appendix}{Appendix III}, 它重印了一篇最初发表于 Birkhäuser 文集的综述文章. 该附录涉及早期版本未讨论的若干方面, 特别包括: 从连续介质力学的基本守恒原理简要导出 Navier--Stokes 方程, 补充一些历史视角, 并介绍若干新进展. 它还概述了不属于本书主题的相关方程, 即 Euler 方程和可压缩 Navier--Stokes 方程. 建议读者在阅读本书主体之前先浏览该附录.

如果现在撰写或重写本书, 就必须面对下述困难: 撰写第一版时, 我们曾在一定程度上试图纳入当时所有关于 Navier--Stokes 方程解的存在性, 唯一性及其逼近的材料. 此后知识体系显著扩展, 如今单独一本书已无法容纳全部内容, 因而必须作出取舍. 如本书其他地方所述, 数值方面已经发展成独立领域, 即计算流体力学. 理论方面也出现了大量新进展, 这些内容在 \MBUnresolvedReference{appendix}{Appendix III} 中有所介绍.

这里列举一些与本书关系较密切的进展. 对本书反复使用的若干技术性结果, 后来出现了新的, 更简洁的证明, 例如 \MBUnresolvedReference{source-note}{Proposition 1.1.1 前的脚注}, \MBUnresolvedReference{remark}{Remark 1.2.7} 和 \MBUnresolvedReference{remark}{Remark 2.1.6(iii)}. 空间周期情形得到了广泛研究: 这一情形在概念上更简单, 并且可以使用 Fourier 级数, 但许多困难与本书研究的无滑移情形相同. 主要简化来自边界层相关困难的消失; 边界层是当时正在发展的另一个主题, 本书没有涉及. 时间与空间解析性方面也取得了新结果, 即时间解析性和空间 Gevrey 正则性. 尽管撰写本书时已经有解析性结果, 新结果的证明与本书的精神更为接近. Navier--Stokes 方程解的大时间行为及其与湍流理论的关系也取得了重大进展. 本书未展开的多数新结果可见 R. Temam (1995) 的讲义, 以及 C. Foias, O. Manley, R. Rosa 和 R. Temam (2001) 即将出版的著作, 它们可以作为本书的后续读物. 最后, 湍流流动控制也是一个后来变得可处理的研究方向, 本书除 \MBUnresolvedReference{appendix}{Appendix III} 末尾的若干评论以及展示大规模数值模拟结果的两幅图外, 未涉及这一主题.

我非常高兴 American Mathematical Society 决定重新出版本书, 并希望新版能够发挥作用. 我要特别感谢发起这一项目的 Susan Friedlander, 以及卓有成效地管理项目的 Sergei Gelfand. 还要感谢帮助我再次通读本书并提出许多订正和意见的年轻同事: Didier Bresch, Brian Ewald, Olivier Goubet, Changbing Hu, François Jauberteau, Jean-Michel Rakotoson, Jie Shen, Shouhong Wang, Xiaoming Wang 和 Mohammed Ziane.

从本书对其工作的众多引用可以看出, 本书深受我从老师 Jacques-Louis Lions 那里所学知识的影响. 再往前追溯 Navier--Stokes 方程的历史, Jean Leray (1906--1998) 在其理论方面作出了大量开创性工作, 参见 \MBUnresolvedReference{appendix-introduction}{Appendix III 的 Introduction}. 他在其他若干数学领域也作出了重要的开创性贡献. 他的文集于 1999 年出版; 这些文集以及其他材料都确认他是二十世纪最杰出的数学家之一.

我曾有幸在 Paris 的 Collège de France 参加 Jean Leray 的研讨班并作报告. 那间古老而充满历史的 ``Salle 5'', 总会使一名年轻研究者感到谦卑, 并留下难以忘怀的经历. 我撰写本书时还是一名年轻研究者, Jean Leray 曾给予我亲切的支持与关注. 为铭记这份恩情, 我以崇高敬意将这个新版献给他.

\begin{flushright}
September 2000
\end{flushright}

\clearpage
\section*{第三次修订版序言}
\phantomsection
\label{front:preface-third-edition}
\addcontentsline{toc}{section}{第三次修订版序言}

本书出版以来, 出现了大量与 Navier--Stokes 方程数值逼近理论有关的论文. 人们对这些方程日益关注, 部分原因是它们在航空科学, 气象学, 热工水力学, 石油工业和等离子体物理等许多当前重要的科学与工业应用中发挥着重要作用. 另一个原因是计算能力的发展: 新型计算机已经提供了更强的计算能力, 而超级计算机将在不久的将来进一步提升这一能力. 在计算机上以数值方式求解流体动力学问题的过程称为计算流体力学(CFD). 近年来这一领域显著扩展, 目前已有数千名研究者, 大量应用和极其庞大的文献, 而且这种扩展很可能持续下去.

本书位于计算流体力学与数学分析的交界处, 而 CFD 与数学分析有着牢固的联系. 即使仅限于不可压缩流体 Navier--Stokes 方程的理论和数值分析, 这些学科的迅速扩展也使得单卷著作不可能对其作全面介绍. 不过, 我们认为本书研究的基本问题在相当长时间内仍会具有意义, 本书以目前的形式仍然有用. 对希望了解最新进展或更专门内容的读者, 新版修订并更新了参考文献. 新版还在 \MBUnresolvedReference{comments}{修订版补充评论, printed page 381} 中对近年来取得进展的方向作了说明; 由于篇幅所限, 这一说明必然不完整.

\begin{flushright}
Paris, January 1984
\end{flushright}

\clearpage
\section*{前言}
\phantomsection
\label{front:foreword}
\addcontentsline{toc}{section}{前言}

本书推导了若干关于黏性不可压缩流体 Navier--Stokes 方程的理论与数值分析结果. 我们将处理以下问题: 一方面, 介绍线性与非线性情形, 定常与含时情形中关于解的存在性, 唯一性以及少数情况下正则性的已知结果; 另一方面, 通过离散化逼近这些问题, 即对空间变量使用有限差分和有限元方法, 对时间变量使用有限差分和分步法. 我们尽可能充分地研究数值过程的稳定性和收敛性. 讨论不会局限于这些理论方面: 特别地, 我们在附录中详细说明一种方法的程序实现. 本书研究的所有方法实际上都已得到应用, 但无法给出所有方法实际实现的细节. 我们介绍的理论结果(存在性, 唯一性等)只是非常基本的结果, 其中没有新的结果; 不过, 我们尽力给出简单而自足的论述. 能量方法和紧性方法位于我们研究的两类问题的核心, 并在两者之间形成自然联系.

下面更详细地说明本书内容. 我们首先考虑线性化定常情形(\MBUnresolvedReference{chapter}{Chapter 1}), 随后考虑非线性定常情形(\MBUnresolvedReference{chapter}{Chapter 2}), 最后考虑完整的非线性含时情形(\MBUnresolvedReference{chapter}{Chapter 3}). 每一阶段都会引入新的数学工具, 这些工具本身有用, 也为后续步骤作准备.

在 \MBUnresolvedReference{chapter}{Chapter 1} 中, 简要介绍存在性和唯一性结果之后, 我们说明利用各种有限差分和有限元方法逼近 Stokes 问题. 这使我们有机会引入无散向量函数的多种逼近方法, 而这些方法对于 \MBUnresolvedReference{chapter}{Chapter 2} 和 \MBUnresolvedReference{chapter}{Chapter 3} 所研究问题的数值方面也至关重要.

在 \MBUnresolvedReference{chapter}{Chapter 2} 中, 我们引入连续情形和离散情形的紧性结果, 随后把前一章在线性情形下得到的结果推广到非线性情形. 本章最后利用分岔方法和拓扑方法证明定常 Navier--Stokes 方程解的非唯一性. 论述基本上是自足的.

\MBUnresolvedReference{chapter}{Chapter 3} 研究完整的非线性含时情形. 我们首先介绍若干能够代表 Navier--Stokes 方程数学理论当时状态的结果, 即存在性和唯一性定理. 随后简要介绍问题的数值方面, 把 \MBUnresolvedReference{chapter}{Chapter 1} 中讨论的空间变量离散化与时间变量的常用离散化方法结合起来. 稳定性和收敛性问题用能量方法处理. 我们还研究分步法和人工可压缩性方法.

以上简述足以表明, 本书绝不是对这一主题的系统研究. 这里没有涉及 Navier--Stokes 方程的许多方面. 关于存在性和唯一性问题的若干有趣方法, 例如半群, 奇异积分算子和 Riemann 流形方法, 均未收入. 在问题的数值方面, 我们没有考虑粒子方法, 也没有考虑 Los Alamos Laboratory 发展出的相关方法.

此外, 我们把讨论严格限制在 Navier--Stokes 方程. 一系列可用相同方法处理的问题没有包括在内, 湍流和高 Reynolds 数流动等困难问题也没有包括在内.

本书内容曾作为 University of Maryland 1972--73 学年第一学期 Navier--Stokes 方程与非线性偏微分方程专题学年的一部分讲授. University of Maryland 出版的相应讲义构成本书的第一版.

我非常感谢 University of Maryland Department of Mathematics 和 Institute of Fluid Dynamics and Applied Mathematics 的同事们, 他们对这些讲义的形成表现出浓厚兴趣. Arlett Williamson 以及 J. Osborn, J. Sather 和 P. Wolfe 对手稿的准备作出了直接贡献. 我要感谢他们纠正我的一些英语错误, 并提出有益的意见和建议; 这些帮助改进了手稿. Mrs Pelissier, Fortin 和 Thomasset 也提出了有用的意见. 最后, 我要感谢 Maryland 和 Orsay 数学系的秘书们在手稿准备过程中给予的全部帮助.

\begin{flushright}
Roger Temam
\end{flushright}

\clearpage
\pagenumbering{arabic}
